%-------------------------------------------------------------------------------
\section{Lintability Determination}
\label{sec:method-lintability}
\label{sec:lintability}
%-------------------------------------------------------------------------------

\begin{table*}[tb]
\centering
\caption{The four lintability conditions, each a Boolean function of one IR field.}
\label{tab:four-conditions}
\footnotesize
\setlength{\tabcolsep}{4pt}
\renewcommand{\arraystretch}{1.2}
\begin{tabularx}{\textwidth}{@{}l l L{0.36\linewidth} X@{}}
\hline
ID & IR field & Value domain & Condition \\
\hline
$C_1$      & obligation         & MUST, SHALL, REQUIRED, SHOULD, RECOMMENDED, MUST NOT, SHALL NOT, SHOULD NOT, MAY, OPTIONAL & obligation $\notin$ \{MAY, OPTIONAL\} \\
$C_2$      & assertion\_subject & Certificate, CA, CrossArtifact & assertion\_subject $=$ Certificate \\
$C_3$      & enforcement\_phase & Encoding, Validation, Comparison, Processing & enforcement\_phase $=$ Encoding \\
$\neg C_4$ & rule\_category     & encoding\_constraint, comparison, definition, capability, algorithm\_ref, display, \dots & rule\_category $\notin$ \{definition, capability, algorithm\_ref, display, \dots\} \\
\hline
\end{tabularx}
\end{table*}

Before generating code, we decide the lintability of each normative rule, namely
whether it admits such a static lint check. Rather than asking the LLM to make
this judgment directly, we compute it from the IR that the LLM has already
produced. Three observations motivate this design. First, the same four pieces
of information determine the lintability of normative rules in practice, namely
the deontic strength of the obligation, who is obligated, when the obligation
applies, and what type of rule it is. Second, each of the four can be encoded as
a single IR field with a small, closed value set, so the LLM performs a
lower-entropy selection of each field value from the enumerated Value domain
column of Table~\ref{tab:four-conditions}. The lintability
decision then becomes a Boolean function of four discrete fields. Third,
separating the four conditions makes any wrong label traceable to exactly one IR
field rather than to an opaque end-to-end LLM call.

Table~\ref{tab:four-conditions} lists the four IR fields, their value domains, and the conditions they decide. The deontic level obligation comes from the IR core four-tuple in Table~\ref{tab:ir-core}. The remaining three fields, assertion\_\allowbreak subject, enforcement\_\allowbreak phase, and rule\_\allowbreak category, are key extension fields. From the enumeration of rule\_\allowbreak category, we designate
\begin{gather*}
\mathcal{N} = \{definition,\,capability,\,algorithm\_ref,\,display,\,\dots\}
\end{gather*}
as the subset whose rule type carries no static certificate lint check on the certificate. This subset covers term definitions, CA capability statements, delegations to external algorithm specifications, UI rendering rules, and similar non-encoding rule types.

The four conditions in Table~\ref{tab:four-conditions} are necessary for lintability, and each is a Boolean function of exactly one IR field:
\begin{equation}\label{eq:lintable}
\mathrm{lintable}(r) \;\Longrightarrow\; C_1(r)\wedge C_2(r)\wedge C_3(r)\wedge \neg C_4(r),
\end{equation}
\begin{align*}
C_1(r) &\equiv \mathrm{is\_normative}(r.\text{obligation}), \\
C_2(r) &\equiv r.\text{assertion\_subject} = Certificate, \\
C_3(r) &\equiv r.\text{enforcement\_phase} = Encoding, \\
C_4(r) &\equiv r.\text{rule\_category} \in \mathcal{N}.
\end{align*}
The obligation field also fixes the lint severity. MUST-class obligations map to Error, and SHOULD-class obligations map to Warning. Severity is therefore not a second classifier, but a direct read of obligation.

The three non-deontic conditions encode three orthogonal boundaries. $C_2$ is the subject boundary, separating obligations on the certificate from obligations on the CA, relying party, or surrounding ecosystem. $C_3$ is the runtime boundary, separating obligations observable in the issued bytes from those that fire during chain processing, name comparison, revocation handling, or CAA retrieval. $\neg C_4$ is the process boundary, separating state assertions from terms in $\mathcal{N}$.

Since each $C_i$ is a deterministic function of one IR field, the lintability
decision---which adopts the conjunction $C_1\wedge C_2\wedge C_3\wedge\neg C_4$ of
these necessary conditions as its operational criterion---is computed from
the IR rather than predicted by the LLM. The same normative rule IR yields the same label bit-for-bit on every run, and any misclassification is traceable to a specific Layer~2 field assignment.
